\centerline{\bf \Huge Вокруг целых чисел}

\problem{\textbf{1}} Сумма 10 различных натуральных чисел равна 148,
причем известно, что наименьшее из этих чисел составляет не менее половины от наибольшего из них. Найдите эти числа.\\

\problem{\textbf{2}} Джеку Воробью нужно было разложить 150 пиастров по 10 кошелькам. После того как он положил некоторое количество пиастров в первый
кошелёк, в каждый следующий он клал больше, чем в предыдущий. В результате оказалось,
что количество пиастров в первом кошельке не меньше, чем половина количества пиастров
в последнем. Сколько пиастров находится в 6-м кошельке?\\

\problem{\textbf{3}} Доля отличников в классе составляет менее
одной трети, но более 20\%, причем известно, что оно кратно 5. А доля хорошистов составляет
менее 20\%, но более 2/11 учащихся. Для школьников, которые не являются хорошистами и
отличниками, организованы дополнительные занятия. Сколько человек должны посещать эти
занятия, если известно, что в классе учится не более 20 школьников?\\


\problem{\textbf{4}} На день рождения Андрея последней пришла Яна, подарившая
ему мяч, а предпоследним — Эдуард, подаривший ему калькулятор. Испытывая калькулятор,
Андрей заметил, что произведение количества всех его подарков на количество подарков, которые были у него до прихода Эдуарда, ровно на 16 больше, чем произведение его возраста на
количество подарков, которые были у него до прихода Яны. Сколько подарков у Андрея?\\

\problem{\textbf{5}} На клетчатой бумаге (сторона клетки 1 см) нарисован прямоугольник, стороны которого лежат на линиях сетки, причём одна сторона на 5 см меньше другой. Оказалось, что его можно разрезать по линиям сетки на несколько частей и сложить из них квадрат. Чему может быть равна сторона этого квадрата? Найдите все возможные значения.

\problem{\textbf{6}}  На доске записано несколько последовательных натуральных
чисел. Ровно 52\% из них — чётные. Сколько чётных чисел записано на доске?

\problem{\textbf{7}}  В школе учится не менее 150 мальчиков, а девочек —
на 15\% больше, чем мальчиков. Когда мальчики поехали на сборы, потребовалось 6 автобусов,
причём в каждом автобусе ехало одинаковое количество школьников. Сколько всего человек
учится в школе, если известно, что общее число учащихся не больше 400?\\


\problem{\textbf{8}}  На доске записано несколько последовательных натуральных
чисел. Ровно 52\% из них — чётные. Сколько чётных чисел записано на доске?\\

\problem{\textbf{9}} Ученикам 11 «Б» класса на выбор предложили пройти тестирование ровно по одному из предметов: биологии, математике или химии. Двое ребят приняли участие в
тестировании по биологии; более трети, но менее 40\% учеников проходили тестирование по химии и ровно половина — по математике. Сколько ребят участвовало в тестировании по химии,
если в классе присутствовали более 16 учеников?
\\

\problem{\textbf{10}} Ученику дано число $x$: это обыкновенная дробь со знаменателем 9. Ученик вычислил три новых числа $2x, 4x$ и $5x$, каждое из этих трёх чисел округлил
до ближайшего целого и результаты округлений сложил. Получилось 111. Найдите $x$. (Число
округляется в меньшую сторону, если его дробная часть меньше 1/2, и в большую, если дробная
часть больше либо равна 1/2.)\\

\problem{\textbf{11}} 779 новогодних подарков разложены по мешкам. В некоторых мешках
лежит по $n$ подарков, в других — по 10 подарков. Какое наименьшее значение может принимать $n$, если всего 25 мешков?\\

\problem{\textbf{12}} Команда спортсменов, третья часть которых — сноубордисты,
спустилась с горы. При этом некоторые из них сели в вагон фуникулёра, вмещающий не более
11 человек, а все остальные спустились самостоятельно, причём их число оказалось больше
40\%, но меньше 44\% от общего количества. Определите количество сноубордистов (если оно
определяется из условия задачи неоднозначно, то впишите в ответ сумму всех возможных его
значений).

